\documentclass[11pt]{scrartcl}

\usepackage[T1]{fontenc}
\usepackage[utf8]{inputenc}
\usepackage{lmodern}
\usepackage{microtype}
\usepackage[margin=1in]{geometry}
\usepackage{setspace}

\title{Teaching Statement}
\author{Decory Edwards}
\date{March 4, 2026}

\setlength{\parindent}{1.5em}
\setlength{\parskip}{0pt}

\begin{document}

\maketitle
\thispagestyle{plain}

\section{Teaching Experience}

As a graduate student, I've been a teaching assistant (TA) mainly for macroeconomics courses. These courses include introductory and intermediate level macroeconomics, as well as a seminar course in macroeconomic models and policy.

I was a teaching assistant for Elements of Macroeconomics for three semesters. The textbook for the course was ``Macroeconomics,'' by R. G. Hubbard. Each semester, the class met twice a week with the professor for lectures. Every Friday, I would lead my own section for the course with students. Additionally, I held an office hour for students each week which lasted around an hour to an hour and a half. The Elements of Macroeconomics course generally had anywhere from 30 to 40 students enrolled each semester.

I also was a teaching assistant for Macroeconomic Theory, which is the next course in the sequence for undergraduates, for three semesters as well. The textbook for the course each semester was ``Macroeconomics,'' 11th edition, by N. Gregory Mankiw. Again, the class met twice a week for lectures and each Friday I led an hour-long section meeting with students. I also held an hour-long office hour session for students each week. The Macroeconomic Theory course had anywhere from 25 to 40 students enrolled each semester.

I was a teaching assistant for Debates in Macroeconomics, an upper-level seminar course for undergraduates. I held an hour-long, weekly office hour and was responsible for helping with grading and proctoring exams. There were 16 students enrolled in this course. In this course, students were challenged to use models from both micro and macro economics to discuss the role for policy in a number of scenarios, such as climate change, immigration, and inequality.

I've also tutored introductory and intermediate courses in mathematics and economics during my undergraduate years. During graduate school, I've done private tutoring as an independent contractor. These included a role as a research mentor for high school students, as well as weekly tutoring for AP and college level macroeconomics, microeconomics, and calculus at various levels.

These experiences not only kept me up to speed in my own studies, but also exposed me to the various opportunities to teach the concepts I've learned over the years. In my view, being able to support other budding learners is my way of paying it forward since I've received similar help in my own studies.

\section{Teaching Philosophy}

I believe teaching is an individualized journey for each student. It is my belief that my job is centered around aiding a student convince themselves that they have transitioned from a state of confusion to understanding. As a teacher, I understand that the most helpful that I can be to a learner is supporting them, encouraging them, and making sure that I've convinced myself of the same material they are trying to learn. These are all necessary to establish the proper trust between teacher and student. In another role as a private tutor, I also learned that many of these concepts are fresh on my mind, and that the ones that aren't can be recovered swiftly. This has led me to seriously consider pursuing opportunities to be an instructor in the future.

\end{document}
